\section{Moving a key between devices}
\label{sec:device-transfer}

\code{keyquorum transfer} copies or moves an active identity between two
devices that are both open. Each side has its own container and its own
database; the global \flag{--db} is not used. Enroll a slot before the
first transfer so this device has an active identity for that label.

\begin{lstlisting}
keyquorum transfer enroll --device ./usb --db ./alice.sqlite --label M.S
keyquorum transfer copy \
  --from-device ./usb --from-db ./alice.sqlite \
  --to-device ./usb2 --to-db ./bob.sqlite \
  --label M.S
keyquorum transfer move \
  --from-device ./usb --from-db ./alice.sqlite \
  --to-device ./usb2 --to-db ./bob.sqlite \
  --label M.S --descendants key-only
keyquorum transfer list --db ./bob.sqlite
keyquorum transfer list --db ./alice.sqlite --all
keyquorum transfer recover \
  --from-device ./usb --from-db ./alice.sqlite \
  --to-device ./usb2 --to-db ./bob.sqlite \
  --transaction <hex-id>
\end{lstlisting}

The identity id is stable; the dotted label is only the current place in the
hierarchy. Possession on a device is \code{active} (the slot secret is
here), \code{ghost} (the hierarchy row remains and the secret does not), or
absent. \code{transfer list} hides ghosts unless \flag{--all} is set.

\begin{itemize}
  \item \textbf{COPY} leaves the source active and records that the
  destination may hold a legitimate copy.
  \item \textbf{MOVE} leaves the source row active until the destination has
  committed and the source slot token is gone. The ghost row is written
  only after that deletion, so a reported ghost cannot still be opened by
  \code{keyquorum-device}. A transfer that stops before the destination
  commits keeps the usable copy on the source.
\end{itemize}

\code{transfer recover} finishes a transfer whose destination already
committed, or rolls back one that did not.

\flag{--descendants} selects what travels with the label: \code{key-only}
(the default), \code{direct-children}, or \code{all-descendants} (the full
subtree). A parent can move without its children, and a child can move
while the parent stays. \flag{--as} names the active identity authorizing
the transfer and defaults to \flag{--label}.

\begin{noteBox}
A ghost keeps hierarchy and provenance and holds no usable private key. It
cannot sign, satisfy a quorum, authorize an import, or be exported. An
active child under a ghost ancestor stays usable. Ghosts do not authorize an
incoming key.
\end{noteBox}

An empty receiver accepts the package. A receiver that already holds active
identities accepts an incoming key only when it is a descendant of one of
those identities. The same identity id with the same public keys reconciles;
the same id or the same label with different key material is refused.

\subsection{Package format and audit}

The package magic is \code{KQTX}. The source device key signs it, and it is
bound to the destination device id. Replay, tamper, and a bad source are
refused. Every attempt is audited without private key material; SQLite
stores the transfer state and the audit row. The raw \code{KQTX} package
stays off the mailbox.

\subsection{Transferring through the mailbox}

When the two devices are not open together, the commands below seal that
package, or a slot relocate, into a \file{.kqpb} letter, and the mailbox
stores only that letter.

\begin{lstlisting}
keyquorum device publish ./usb --url https://relay.example.com
keyquorum transfer relay-send \
  --operation move \
  --from-device ./usb --from-db ./alice.sqlite \
  --label M.S \
  --to-device-id <hex-device-id> \
  --recipient-key-file ./bob-slot.pub \
  --url https://relay.example.com
keyquorum transfer relay-collect \
  --to-device ./usb2 --to-db ./bob.sqlite \
  --slot recv \
  --from-device-id <hex-source-id> \
  --url https://relay.example.com
keyquorum transfer relay-finalize \
  --from-device ./usb --from-db ./alice.sqlite \
  --slot M.S \
  --to-device-id <hex-device-id> \
  --url https://relay.example.com
keyquorum device relay-relocate \
  --from ./usb --label M.S \
  --to-device-id <hex-device-id> \
  --recipient-key-file ./bob-slot.pub \
  --url https://relay.example.com
keyquorum device relay-accept ./usb2 --slot recv \
  --url https://relay.example.com
keyquorum device relay-drop ./usb --label M.S \
  --to-device-id <hex-device-id> \
  --relocate-id <hex-relocate-id> \
  --url https://relay.example.com
\end{lstlisting}

\code{relay-send} leaves the source prepared. COPY stays active. MOVE
deletes the source slot token and writes the ghost only when
\code{relay-finalize} checks the destination's signed acknowledgement.
\code{relay-relocate} prints a relocate id and leaves the source slot in
place until \code{relay-drop} checks an acknowledgement signed for that id
--- an acknowledgement from an earlier relocation of the same slot does not
match. \code{relay-accept} acknowledges a slot that is already installed
with the same keys again, so a retry after a failed upload finishes.

The mailbox keeps device letters and acknowledgements for 30 days and never
deletes them on acknowledgement. Load a \code{device.push} key for sends and
publishes, and a \code{device.pull} key bound to the recipient fingerprint
for collects, finalizes, and drops. \code{relay-collect} and
\code{relay-accept} also need \code{device.push} to post the
acknowledgement (\flag{--push-key}, or a stored key of that scope).
